\begin{table}[htbp]
\centering
\caption{Standardized Effect Sizes: Sheppard-Towner Act Exposure}
\label{tab:sde}
\begin{threeparttable}
\begin{adjustbox}{max width=\textwidth}
\begin{tabular}{lcccccc}
\toprule
Outcome & $\hat{\beta}$ & SE & SD($Y$) & SDE & SE(SDE) & Classification \\
\midrule
\multicolumn{7}{l}{\textit{Panel A: Pooled}} \\
Wage Income & 38.6 & 17.0 & 1584.4 & 0.0244 & 0.0108 & Small positive \\
Education (years) & -0.068 & 0.168 & 4.48 & -0.0152 & 0.0374 & Small negative \\
Occ. Income Score & 0.06 & 0.10 & 9.43 & 0.0067 & 0.0105 & Small positive \\
\midrule
\multicolumn{7}{l}{\textit{Panel B: Heterogeneous (Wage Income)}} \\
Black subsample & 91.9 & 41.3 & 1368.5 & 0.0671 & 0.0302 & Moderate positive \\
Rural-born subsample & 70.4 & 17.7 & 1578.1 & 0.0446 & 0.0112 & Small positive \\
\bottomrule
\end{tabular}
\end{adjustbox}
\begin{tablenotes}[flushleft]\footnotesize
\item \textit{Notes:} \textbf{Country:} United States. \textbf{Research question:} Does prenatal and infant exposure to America's first federal public health program (Sheppard-Towner Act, 1921--1929) produce lasting human capital gains measurable in adult wage income, education, and occupational attainment? \textbf{Policy mechanism:} The Act provided federal matching funds to states for maternal and infant health education, establishing approximately 3,000 prenatal clinics and deploying visiting nurses for over 3 million home visits targeting expectant and new mothers. Three states (Massachusetts, Connecticut, Illinois) refused participation on states' rights grounds. \textbf{Outcome definition:} Panel A reports wage income in 1950 dollars (incwage\_1950), years of completed schooling (educ\_1950), and occupational income score (occscore\_1950). Panel B reports wage income for Black individuals and rural-born individuals separately. \textbf{Treatment:} Binary --- born 1922--1928 in a state that accepted Sheppard-Towner funds vs.\ born in a non-participating state or in a pre-treatment cohort. \textbf{Data:} IPUMS Multigenerational Longitudinal Panel (MLP) linking individuals across 1930, 1940, and 1950 full-count censuses; birth cohorts 1912--1928; US-born individuals. \textbf{Method:} Triple-differences (DDD) with birth-year and birth-state fixed effects, demographic controls, standard errors clustered by birth state. \textbf{Sample:} US-born individuals observed in all three census decades, restricted to birth cohorts 1912--1928, with valid wage income in 1950. SDE $= \hat{\beta} / \text{SD}(Y)$ where SD($Y$) is the pre-treatment standard deviation (non-participating states, pre-1922 cohorts). Classification refers to magnitude, not statistical significance: Large ($|$SDE$| > 0.15$), Moderate ($0.05$--$0.15$), Small ($0.005$--$0.05$), Null ($< 0.005$). 
\end{tablenotes}
\end{threeparttable}
\end{table}
